% ============================================================================
%  Chapitre 5, partie III : Mise en équation linéaire
% ============================================================================

% ----------------------------------------------------------------------------
\section{La mise en équations linéaires des contrôles}
% ----------------------------------------------------------------------------

On écrit maintenant les deux contrôles du modèle,
$\vect{L}_G\cdot\vect{e}_z$ et $T$, avec les seules mesures
$\vect{\omega}$ et $\vect{\gamma}$. Dans le repère du gyromètre,
$\mat{I}_b$ est remplacée par $\mat{J}$, et les deux contrôles
deviennent linéaires en ses coefficients.

\subsection{Moment cinétique vertical}

D'après la partie précédente, avec $\mat{J}$ à la place de
$\mat{I}_b$, dans le repère gyromètre
\[
    \vect{L}_G\cdot\vect{e}_z = \vect{\gamma}^T\mat{J}\,\vect{\omega} = c
\]
où $c$ est une constante inconnue. $\mat{J}$ n'est pas diagonale,
mais elle est symétrique donc 
\[
    J_{11}\gamma_1\omega_1 + J_{22}\gamma_2\omega_2 + J_{33}\gamma_3\omega_3
    + J_{12}(\gamma_1\omega_2 + \gamma_2\omega_1)
    + J_{23}(\gamma_2\omega_3 + \gamma_3\omega_2)
    + J_{13}(\gamma_1\omega_3 + \gamma_3\omega_1) - c = 0
\]

\subsection{Énergie cinétique}

\textbf{Vitesse du centre : } la boule roule sans glisser, donc
$\vect{v}_P = 0$, et $\overrightarrow{PG} = a\,\vect{e}_z$ donne
\[
    \vect{v}_G = \vect{v}_P + \vect{\Omega}\wedge\overrightarrow{PG}
    = a\,\vect{\Omega}\wedge\vect{e}_z
\]
Avec $\vect{\Omega} = \mat{R}\vect{\omega}$ et
$\vect{e}_z = \mat{R}\vect{\gamma}$, on a 
\[
    \vect{v}_G = a\,(\mat{R}\vect{\omega})\wedge(\mat{R}\vect{\gamma})
    = a\,\mat{R}\,(\vect{\omega}\wedge\vect{\gamma})
\]
Or $\mat{R}^T\mat{R} = \Id$, donc
\[
    \|\vect{v}_G\|^2 = a^2\|\vect{\omega}\wedge\vect{\gamma}\|^2
\]
et la relation
$\|\vect{u}\wedge\vect{v}\|^2 = \|\vect{u}\|^2\|\vect{v}\|^2 - (\vect{u}\cdot\vect{v})^2$,
avec $\|\vect{\gamma}\| = 1$, donne
\[
    \|\vect{v}_G\|^2 = a^2\bigl(\|\vect{\omega}\|^2 - (\vect{\gamma}^T\vect{\omega})^2\bigr)
\]

\textbf{Énergie : } l'énergie cinétique, conservée, s'écrit alors
\[
    2T = m\|\vect{v}_G\|^2 + \vect{\omega}^T\mat{J}\,\vect{\omega}
    = m a^2\bigl(\|\vect{\omega}\|^2 - (\vect{\gamma}^T\vect{\omega})^2\bigr)
    + \vect{\omega}^T\mat{J}\,\vect{\omega}
\]
En développant $\vect{\omega}^T\mat{J}\,\vect{\omega}$ et en passant
le terme connu à droite,
\[
    J_{11}\omega_1^2 + J_{22}\omega_2^2 + J_{33}\omega_3^2
    + 2J_{12}\omega_1\omega_2 + 2J_{23}\omega_2\omega_3
    + 2J_{13}\omega_1\omega_3 - 2T
    = -m a^2\bigl(\|\vect{\omega}\|^2 - (\vect{\gamma}^T\vect{\omega})^2\bigr)
\]

\subsection{Vecteur des paramètres}

Les deux équations sont linéaires en
\[
    \vect{\theta} = \begin{pmatrix}
        J_{11} & J_{22} & J_{33} & J_{12} & J_{23} & J_{13} & c & T
    \end{pmatrix}^T \in \R^8
\]
où les constantes $c$ et $T$, inconnues, sont estimées en même temps
que $\mat{J}$. Pour chaque mesure, on a donc
\[
    \vect{\varphi}_1(\vect{\gamma}, \vect{\omega})^T\vect{\theta} = 0,
    \qquad
    \vect{\varphi}_2(\vect{\gamma}, \vect{\omega})^T\vect{\theta} = y_2(\vect{\gamma}, \vect{\omega})
\]
avec
\begin{align*}
    \vect{\varphi}_1 &= \bigl(\gamma_1\omega_1,\ \gamma_2\omega_2,\
        \gamma_3\omega_3,\ \gamma_1\omega_2 + \gamma_2\omega_1,\
        \gamma_2\omega_3 + \gamma_3\omega_2,\
        \gamma_1\omega_3 + \gamma_3\omega_1,\ -1,\ 0\bigr)^T\\
    \vect{\varphi}_2 &= \bigl(\omega_1^2,\ \omega_2^2,\ \omega_3^2,\
        2\omega_1\omega_2,\ 2\omega_2\omega_3,\ 2\omega_1\omega_3,\
        0,\ -2\bigr)^T\\
    y_2 &= -m a^2\bigl(\|\vect{\omega}\|^2 - (\vect{\gamma}^T\vect{\omega})^2\bigr)
\end{align*}
